\documentclass[journal,10pt,twocolumn]{IEEEtran}

\usepackage[utf8]{inputenc}
\usepackage[spanish,es-tabla]{babel}
\usepackage{amsmath,amssymb,amsfonts}
\usepackage{algorithmic}
\usepackage{graphicx}
\usepackage{textcomp}
\usepackage{xcolor}
\usepackage{booktabs}
\usepackage{cite}
\usepackage{url}
\usepackage{microtype}

\begin{document}

\title{Data-Efficient Contrastive Learning: Evaluación de la Reducción de Datos mediante SAS en un Nuevo Dominio y Múltiples Paradigmas Auto-Supervisados}

\author{
    \IEEEauthorblockN{Manuel Fernando Rocha Zuleta\IEEEauthorrefmark{1}, Juan Pablo Hoyos Sánchez\IEEEauthorrefmark{2}} \\
    \IEEEauthorblockA{\textit{Programa Curricular de Ingeniería Mecatrónica} \\
    \textit{Facultad de Minas / Ingeniería, Universidad Nacional de Colombia} \\
    Medellín, Colombia \\
    \IEEEauthorrefmark{1}mfrochaz@unal.edu.co, \IEEEauthorrefmark{2}jphoyoss@unal.edu.co}
}

\maketitle

\begin{abstract}
El aprendizaje auto-supervisado (SSL) permite entrenar modelos profundos de visión artificial sin requerir anotaciones manuales, pero impone una carga computacional masiva debido a la necesidad de procesar múltiples vistas aumentadas sobre conjuntos de datos a gran escala. En aplicaciones mecatrónicas y de visión embebida, el costo de cómputo y el tiempo de convergencia son cuellos de botella críticos. Este trabajo investiga el impacto de la reducción sistemática de datos mediante el algoritmo SAS (\textit{Subsets that maximize Augmentation Similarity}) en cuatro familias de arquitecturas SSL representativas: SimSiam (no contrastivo siamés), BYOL (basado en bootstrap con red target), CPC (predictivo contrastivo espacial) y Align-Uniform (optimización directa en la hiperesfera). Empleando el conjunto de datos Food-101 (75\,750 imágenes de entrenamiento), se evalúa la retención de representaciones transferibles al reducir el volumen de datos en 20\% (80\% conservado) y 40\% (60\% conservado), en comparación con el conjunto completo y con controles de submuestreo aleatorio uniforme. La calidad de las representaciones se evalúa mediante el protocolo estándar de \textit{linear probing}. Los resultados experimentales y la metodología reproducible se analizan desde la perspectiva del diseño de sistemas de percepción eficientes para visión computacional.
\end{abstract}

\begin{IEEEkeywords}
Aprendizaje auto-supervisado (SSL), reducción de datos, SAS, selección de subconjuntos, visión artificial, Food-101, linear probing, eficiencia computacional.
\end{IEEEkeywords}

\section{Introducción}
\label{sec:intro}

En la ingeniería mecatrónica moderna, la percepción visual autónoma es un componente crítico para sistemas de robótica móvil, inspección automatizada en manufactura y clasificación de productos. Tradicionalmente, los modelos de visión por computador han dependido del aprendizaje supervisado, el cual exige grandes volúmenes de datos anotados manualmente por expertos humanos. Este proceso de etiquetado manual introduce elevados costos temporales y económicos, limitando la adaptabilidad de los sistemas inteligentes ante nuevos entornos o dominios visuales.

El Aprendizaje Auto-Supervisado (\textit{Self-Supervised Learning}, SSL) ha emergido como una alternativa fundamental, capaz de aprender representaciones visuales invariantes y discriminativas a partir de datos no etiquetados \cite{chen2021simsiam, grill2020byol}. No obstante, el entrenamiento de modelos SSL sobre conjuntos de datos a gran escala conlleva un elevado costo computacional, energético y económico, restringiendo su escalabilidad y sostenibilidad práctica \cite{mirzasoleiman2020coresets}.

Recientes avances teóricos han demostrado que no todos los ejemplos de entrenamiento contribuyen equitativamente al aprendizaje de representaciones. En particular, Joshi y Mirzasoleiman \cite{joshi2023sas} probaron que los ejemplos de mayor beneficio para el aprendizaje contrastivo son aquellos cuya similitud esperada entre vistas aumentadas es máxima respecto al resto de la distribución. Sorprendentemente, estos ejemplos coinciden con los catalogados como ``fáciles'' en aprendizaje supervisado, permitiendo excluir hasta el 40\% de los datos de entrenamiento sin degradar el rendimiento en tareas posteriores (\textit{downstream}).

Sin embargo, los resultados originales de Joshi y Mirzasoleiman \cite{joshi2023sas} se limitaron a conjuntos sintéticos o de baja resolución (CIFAR-100, STL-10 y TinyImageNet) y al algoritmo SimCLR. Queda sin verificar empíricamente si las propiedades de la similitud de aumentos se generalizan a dominios con alta variabilidad intra-clase y granularidad fina, así como su independencia arquitectónica frente a otros paradigmas modernos de aprendizaje auto-supervisado como SimSiam \cite{chen2021simsiam}, BYOL \cite{grill2020byol}, Contrastive Predictive Coding (CPC) \cite{oord2018cpc} y la optimización directa de alineación y uniformidad en la hiperesfera \cite{wang2020alignuniform}.

Este trabajo de investigación formaliza la evaluación de SAS sobre el conjunto de datos Food-101 \cite{bossard2014food101}, persiguiendo los siguientes objetivos específicos:
\begin{enumerate}
    \item Implementar el algoritmo SAS en modo puramente no supervisado y generar subconjuntos de datos con reducciones del 20\% y 40\% (80\% y 60\% de datos retenidos) sobre Food-101.
    \item Entrenar cuatro familias representativas de modelos auto-supervisados (SimSiam, BYOL, CPC y Align-Uniform) tanto en el conjunto completo como en los subconjuntos optimizados por SAS y subconjuntos aleatorios uniformes de control.
    \item Evaluar y comparar cuantitativamente el rendimiento downstream mediante el protocolo estándar de evaluación lineal (\textit{linear probing}), analizando la retención de precisión Top-1 / Top-5 y las ganancias en eficiencia computacional.
\end{enumerate}

\section{Trabajo Relacionado}
\label{sec:related_work}

\subsection{Paradigmas en Aprendizaje Auto-Supervisado (SSL)}
Los métodos modernos de SSL para visión por computador pueden categorizarse principalmente según la estrategia empleada para prevenir el colapso de representación (donde la red proyecta todas las entradas a una constante trivial):

\subsubsection{Métodos Siameses y de Detención de Gradiente}
Chen y He \cite{chen2021simsiam} introdujeron SimSiam, demostrando que una arquitectura siamesa estándar puede aprender representaciones no colapsadas sin necesidad de pares negativos, capas de regularización complejas ni redes de memoria, recurriendo exclusivamente a un proyector asimétrico y la operación de detención de gradiente (\textit{stop-gradient}).

\subsubsection{Métodos de Red Objetivo Asimétrica}
Grill et al. \cite{grill2020byol} propusieron BYOL (\textit{Bootstrap Your Own Latent}), donde una red \textit{online} aprende a predecir la representación producida por una red \textit{target}. Los parámetros de la red objetivo $\theta_{\text{target}}$ se actualizan como un promedio móvil exponencial (EMA) de los pesos de la red en línea $\theta_{\text{online}}$:
\begin{equation}
\theta_{\text{target}} \leftarrow \tau \theta_{\text{target}} + (1 - \tau) \theta_{\text{online}}
\label{eq:byol_ema}
\end{equation}
donde $\tau \in [0, 1]$ es el coeficiente de momento.

\subsubsection{Codificación Predictiva Contrastiva (CPC)}
Oord et al. \cite{oord2018cpc} formalizaron CPC, fundamentado en predecir el futuro en representaciones latentes secuenciales o la correlación espacial en imágenes divididas en franjas o parches, maximizando una cota inferior de la información mutua mediante la pérdida InfoNCE.

\subsubsection{Alineamiento y Uniformidad en la Hiperesfera}
Wang e Isola \cite{wang2020alignuniform} demostraron teóricamente que las funciones de pérdida contrastivas optimizan dos propiedades geométricas asintóticas en la hiperesfera $\mathcal{S}^{d-1}$:
\begin{itemize}
    \item \textbf{Alineamiento} (\textit{Alignment}): dos vistas aumentadas del mismo objeto deben mapearse a representaciones cercanas en la esfera.
    \item \textbf{Uniformidad} (\textit{Uniformity}): las representaciones de muestras distintas deben preservar la máxima información distribuyéndose uniformemente sobre la esfera.
\end{itemize}

\subsection{Selección de Subconjuntos y Coreset Selection en SSL}
La mayoría de algoritmos clásicos de selección de datos (p. ej., basados en pérdida o gradientes) dependen críticamente de etiquetas supervisadas. Para escenarios puramente no supervisados, Joshi y Mirzasoleiman \cite{joshi2023sas} demostraron que el submuestreo aleatorio descarta muestras con alta capacidad de aprendizaje y retiene redundancias. SAS aborda este problema formulando la selección de muestras como un problema de importancia muestral condicionado a clases latentes.

\section{Metodología y Formulación Matemática}
\label{sec:methodology}

\subsection{Arquitectura del Sistema de Visión}
El sistema emplea una red troncal compartida (\textit{backbone}) ResNet-18 \cite{he2016resnet} sin la capa lineal final de clasificación, produciendo un vector de características latentes $h \in \mathbb{R}^{512}$. Sobre esta base se integran cabezas de proyección específicas para cada técnica.

\subsection{Selección de Subconjuntos mediante SAS No Supervisado}
Dado un conjunto no etiquetado $\mathcal{D} = \{x_i\}_{i=1}^N$ con $N = 75\,750$ imágenes de entrenamiento, el algoritmo SAS \cite{joshi2023sas} en su modo puramente auto-supervisado consta de tres etapas:

\begin{enumerate}
    \item \textbf{Extracción con Modelo Proxy}: Se emplea una red ResNet-18 preentrenada en ImageNet como extractor proxy para mapear cada imagen $x_i$ a una representación $z_i \in \mathbb{R}^{512}$, normalizada en la esfera unitaria: $z_i \leftarrow z_i / \|z_i\|_2$.
    \item \textbf{Aproximación de Clases Latentes}: Sin acceder a las etiquetas ground-truth de Food-101, se particiona el espacio latente en $K = 101$ conglomerados mediante el algoritmo $K$-Means:
    \begin{equation}
    \min_{\{C_k\}_{k=1}^K} \sum_{k=1}^K \sum_{z_i \in C_k} \|z_i - \mu_k\|_2^2
    \end{equation}
    donde $\mu_k$ representa el centroide del grupo $k$.
    \item \textbf{Muestreo por Importancia Intra-Cluster}: Para cada conglomerado $C_k$, se asigna una probabilidad de selección proporcional a la similitud de aumento, priorizando muestras centrales y representativas que minimizan la varianza del gradiente de la pérdida contrastiva. La cuota por cluster se asigna de forma estratificada hasta alcanzar el presupuesto global $B = \lfloor \alpha N \rfloor$, donde $\alpha \in \{0.80, 0.60\}$.
\end{enumerate}

\subsection{Formulación de los Modelos SSL Evaluados}

\subsubsection{SimSiam}
Dadas dos vistas aumentadas $x_1, x_2$, las proyecciones son $z_1 = g(f(x_1))$, $z_2 = g(f(x_2))$, y las predicciones son $p_1 = h(z_1)$, $p_2 = h(z_2)$. La función de pérdida simétrica es:
\begin{equation}
\mathcal{L}_{\text{SimSiam}} = \frac{1}{2} \mathcal{D}(p_1, \text{stopgrad}(z_2)) + \frac{1}{2} \mathcal{D}(p_2, \text{stopgrad}(z_1))
\end{equation}
donde $\mathcal{D}(p, z) = -\frac{\langle p, z \rangle}{\|p\|_2 \|z\|_2}$.

\subsubsection{BYOL}
Emplea una red \textit{online} $(f_\theta, g_\theta, q_\theta)$ y una red \textit{target} $(f_\xi, g_\xi)$. La función de pérdida normalizada por pares es:
\begin{equation}
\mathcal{L}_{\text{BYOL}} = 2 - 2 \frac{\langle q_\theta(z_1), z'_2 \rangle}{\|q_\theta(z_1)\|_2 \|z'_2\|_2} + 2 - 2 \frac{\langle q_\theta(z_2), z'_1 \rangle}{\|q_\theta(z_2)\|_2 \|z'_1\|_2}
\end{equation}
donde $z'_i = g_\xi(f_\xi(x_i))$ proviene de la red congelada respecto a gradientes directos.

\subsubsection{CPC (Contrastive Predictive Coding)}
La imagen de entrada se particiona verticalmente en $R = 4$ franjas horizontales. Cada franja se codifica en un vector $z_r$. Una red autorregresiva (GRU o MLP causal) resume el contexto superior $c_t$ y predice las representaciones futuras mediante la pérdida InfoNCE:
\begin{equation}
\mathcal{L}_{\text{CPC}} = -\sum_{k} \log \frac{\exp(z_{t+k}^T W_k c_t)}{\exp(z_{t+k}^T W_k c_t) + \sum_{j} \exp(z_j^T W_k c_t)}
\end{equation}

\subsubsection{Align-Uniform}
Optimiza directamente las métricas en $\mathcal{S}^{d-1}$:
\begin{align}
\mathcal{L}_{\text{align}}(f; \alpha) &\triangleq \mathbb{E}_{(x, x^+)} \left[ \|f(x) - f(x^+)\|_2^\alpha \right], \quad \alpha = 2 \\
\mathcal{L}_{\text{uniform}}(f; t) &\triangleq \log \mathbb{E}_{x, y \overset{\text{i.i.d.}}{\sim} p_{\text{data}}} \left[ e^{-t \|f(x) - f(y)\|_2^2} \right], \quad t = 2
\end{align}
La función multiobjetivo resultante es $\mathcal{L}_{\text{AU}} = \mathcal{L}_{\text{align}} + \lambda \mathcal{L}_{\text{uniform}}$.

\section{Montaje Experimental}
\label{sec:setup}

\subsection{Conjunto de Datos y Preprocesamiento}
Se utiliza el conjunto de datos Food-101 \cite{bossard2014food101}, compuesto por 101 clases de alimentos con un total de 101\,000 imágenes (75\,750 para entrenamiento y 25\,250 para prueba oficial). Para el preentrenamiento auto-supervisado, las imágenes se reescalan a $128 \times 128$ píxeles para optimizar el rendimiento computacional. El aumento de datos incluye:
\begin{itemize}
    \item Recorte aleatorio redimensionado (\textit{RandomResizedCrop}, escala $[0.2, 1.0]$).
    \item Volteo horizontal aleatorio ($p = 0.5$).
    \item Distorsión cromática (\textit{ColorJitter}: brillo 0.4, contraste 0.4, saturación 0.4, tono 0.1 con $p = 0.8$).
    \item Conversión a escala de grises aleatoria ($p = 0.2$).
    \item Desenfoque Gaussiano aleatorio ($\sigma \in [0.1, 2.0]$, $p = 0.5$).
    \item Normalización con la media y desviación estándar de ImageNet.
\end{itemize}

\subsection{Plataformas de Cómputo Distribuidas}
Para viabilizar el volumen experimental sin sobrepasar límites de cuota, el entrenamiento se distribuyó en dos estaciones de trabajo heterogéneas:
\begin{enumerate}
    \item \textbf{Estación A (Kaggle Cloud)}: GPU NVIDIA Tesla T4 / P100 (16\,GB VRAM), asignada a los modelos \textit{SimSiam} y \textit{BYOL}.
    \item \textbf{Estación B (Laboratorio de Mecatrónica)}: Estación de trabajo local con GPU NVIDIA RTX A2000 (12\,GB VRAM, arquitectura Ampere) y almacenamiento NVMe de 1\,TB, asignada a los modelos \textit{CPC} y \textit{Align-Uniform}.
\end{enumerate}

\subsection{Hiperparámetros de Entrenamiento}
\begin{table}[ht]
\centering
\caption{Parámetros de configuración para el entrenamiento SSL (\texttt{config\_final.py})}
\label{tab:hyperparams}
\begin{tabular}{ll}
\toprule
\textbf{Parámetro} & \textbf{Valor / Especificación} \\
\midrule
Arquitectura base & ResNet-18 (512 dimensiones latentes) \\
Dimensión oculta proyector & 2048 \\
Dimensión salida proyector & 256 \\
Épocas de preentrenamiento SSL & 200 épocas \\
Tamaño de lote (\textit{batch size}) & 128 \\
Optimizador & Adam ($\beta_1 = 0.9, \beta_2 = 0.999$) \\
Tasa de aprendizaje base & $3 \times 10^{-4}$ \\
Decaimiento de peso (\textit{weight decay}) & $1 \times 10^{-6}$ \\
Precisión mixta automática (AMP) & Habilitada (\texttt{torch.float16}) \\
Semilla aleatoria & 42 \\
\bottomrule
\end{tabular}
\end{table}

\subsection{Protocolo de Evaluación Lineal (Linear Probing)}
Finalizada la etapa de preentrenamiento auto-supervisado, se congela completamente la red troncal ResNet-18 ($\nabla_\theta = 0$) y se acopla una capa lineal de clasificación $\mathbb{R}^{512} \to \mathbb{R}^{101}$. Esta cabeza lineal se entrena durante 100 épocas sobre el conjunto de entrenamiento de Food-101 supervisado, evaluando la exactitud de clasificación (Top-1 y Top-5) sobre las 25\,250 muestras de prueba independientes.

\section{Resultados Experimentales}
\label{sec:results}

\subsection{Matriz de Rendimiento Comparativo}
En la Tabla~\ref{tab:results_matrix} se presenta la estructura consolidada para la comparación de los 16 experimentos. 

\begin{table*}[t]
\centering
\caption{Resultados de evaluación lineal (Top-1 / Top-5 Accuracy) y costo computacional en Food-101. \newline \textit{Nota: Los valores numéricos se encuentran actualmente en cómputo sobre Kaggle y la estación RTX A2000.}}
\label{tab:results_matrix}
\begin{tabular}{llccccc}
\toprule
\textbf{Método SSL} & \textbf{Configuración de Datos} & \textbf{Muestras} & \textbf{Tiempo Preentrenamiento} & \textbf{Pérdida SSL Final} & \textbf{Top-1 Acc (\%)} & \textbf{Top-5 Acc (\%)} \\
\midrule
\textbf{SimSiam} & Completo (100\%) & 75\,750 & \textit{[En ejecución]} & \textit{[En ejecución]} & \textit{[Pendiente]} & \textit{[Pendiente]} \\
                 & SAS 80\% & 60\,600 & \textit{[En ejecución]} & \textit{[En ejecución]} & \textit{[Pendiente]} & \textit{[Pendiente]} \\
                 & SAS 60\% & 45\,450 & \textit{[En ejecución]} & \textit{[En ejecución]} & \textit{[Pendiente]} & \textit{[Pendiente]} \\
                 & Aleatorio 60\% & 45\,450 & \textit{[En ejecución]} & \textit{[En ejecución]} & \textit{[Pendiente]} & \textit{[Pendiente]} \\
\midrule
\textbf{BYOL}    & Completo (100\%) & 75\,750 & \textit{[En ejecución]} & \textit{[En ejecución]} & \textit{[Pendiente]} & \textit{[Pendiente]} \\
                 & SAS 80\% & 60\,600 & \textit{[En ejecución]} & \textit{[En ejecución]} & \textit{[Pendiente]} & \textit{[Pendiente]} \\
                 & SAS 60\% & 45\,450 & \textit{[En ejecución]} & \textit{[En ejecución]} & \textit{[Pendiente]} & \textit{[Pendiente]} \\
                 & Aleatorio 60\% & 45\,450 & \textit{[En ejecución]} & \textit{[En ejecución]} & \textit{[Pendiente]} & \textit{[Pendiente]} \\
\midrule
\textbf{CPC}     & Completo (100\%) & 75\,750 & \textit{[En ejecución]} & \textit{[En ejecución]} & \textit{[Pendiente]} & \textit{[Pendiente]} \\
                 & SAS 80\% & 60\,600 & \textit{[En ejecución]} & \textit{[En ejecución]} & \textit{[Pendiente]} & \textit{[Pendiente]} \\
                 & SAS 60\% & 45\,450 & \textit{[En ejecución]} & \textit{[En ejecución]} & \textit{[Pendiente]} & \textit{[Pendiente]} \\
                 & Aleatorio 60\% & 45\,450 & \textit{[En ejecución]} & \textit{[En ejecución]} & \textit{[Pendiente]} & \textit{[Pendiente]} \\
\midrule
\textbf{Align-Uniform} & Completo (100\%) & 75\,750 & \textit{[En ejecución]} & \textit{[En ejecución]} & \textit{[Pendiente]} & \textit{[Pendiente]} \\
                       & SAS 80\% & 60\,600 & \textit{[En ejecución]} & \textit{[En ejecución]} & \textit{[Pendiente]} & \textit{[Pendiente]} \\
                       & SAS 60\% & 45\,450 & \textit{[En ejecución]} & \textit{[En ejecución]} & \textit{[Pendiente]} & \textit{[Pendiente]} \\
                       & Aleatorio 60\% & 45\,450 & \textit{[En ejecución]} & \textit{[En ejecución]} & \textit{[Pendiente]} & \textit{[Pendiente]} \\
\bottomrule
\end{tabular}
\end{table*}

\subsection{Curvas de Aprendizaje y Pérdida}
\textit{[Sección reservada para la inserción de las figuras generadas a partir de los historiales guardados en \texttt{output/results\_final/history\_*.npy} una vez culminen los procesos de entrenamiento]}.

\section{Discusión y Conclusiones}
\label{sec:discussion}

\subsection{Discusión Técnica y Perspectiva en Ingeniería Mecatrónica}
La reducción del volumen de datos en el entrenamiento de modelos de visión artificial tiene repercusiones directas en el ciclo de vida de desarrollo de sistemas mecatrónicos:
\begin{itemize}
    \item \textbf{Eficiencia Energética y de Cómputo}: La reducción en tiempo de GPU se traduce directamente en menor consumo energético en servidores de entrenamiento y prototipado rápido.
    \item \textbf{Despliegue en Borde (\textit{Edge Computing})}: Aunque la inferencia en tiempo real depende del tamaño del \textit{backbone} (ResNet-18), la capacidad de reentrenar o adaptar modelos periódicamente en plantas industriales o robots de servicio con recursos limitados se vuelve viable cuando se requiere solo entre 60\% y 80\% de los datos.
\end{itemize}

\subsection{Limitaciones del Estudio}
\begin{enumerate}
    \item El análisis se acota a la arquitectura ResNet-18 y resolución de $128 \times 128$; arquitecturas basadas en Vision Transformers (ViT) o resoluciones estándar de $224 \times 224$ no fueron evaluadas debido a restricciones de cuota computacional.
    \item El algoritmo SAS requiere una pasada previa de extracción con un modelo proxy preentrenado (ResNet-18 ImageNet), lo que introduce un costo computacional inicial amortizado a lo largo de los experimentos.
\end{enumerate}

\subsection{Conclusiones}
\textit{[Conclusiones cuantitativas específicas pendientes de la consolidación de las métricas Top-1 / Top-5 y tiempos de GPU en la Tabla~\ref{tab:results_matrix}]}.


\bibliographystyle{IEEEtran}
\bibliography{references}

\end{document}
